% answers_short.tex —— 2008 年·普通高等学校招生全国统一考试（四川延考卷）·理科 · 简短答案速查
% 由 tools/make_short_answers.py 自 answers.tex 的 \gksolhead/\gkanswer 提取，
% 并与 2008/2008_数学真题_四川延考区理_答案要点.md 逐题比对后生成，勿手改（改源 MD 后重跑）。

\begin{gknotice}
\textbf{使用说明}：本节只给最终答案便于快速核对，完整推导见解析版。答案为编写者独立推导，\textbf{非官方标准答案}。
\end{gknotice}

\gkpart{选择题}
\begin{center}\begin{tabular}{ccccccccccccc}
\toprule
题号 & 1 & 2 & 3 & 4 & 5 & 6 & 7 & 8 & 9 & 10 & 11 & 12 \\
\midrule
答案 & B & D & C & B & C & A & A & D & B & C & B & D \\
\bottomrule
\end{tabular}\end{center}

\gkpart{填空题}
\begin{itemize}[leftmargin=2.6em]
  \item[\textbf{13.}] $y=\ln(x+1)-1$（$x>-1$）。
  \item[\textbf{14.}] $3$。
  \item[\textbf{15.}] $\dfrac12$。
  \item[\textbf{16.}] $\dfrac{\sqrt2}2$。
\end{itemize}

\gkpart{解答题}
\begin{itemize}[leftmargin=2.6em]
  \item[\textbf{17.}] （Ⅰ）$A=\dfrac{5\pi}8$，$C=\dfrac{\pi}8$；（Ⅱ）$\triangle ABC$ 面积的最大值为 $\sqrt3$。
  \item[\textbf{18.}] （Ⅰ）$0.9$；（Ⅱ）$\xi\sim B(3,0.1)$，分布列 $P(\xi=0)=0.729$、$P(\xi=1)=0.243$、$P(\xi=2)=0.027$、$P(\xi=3)=0.001$，$E\xi=0.3$。
  \item[\textbf{19.}] （Ⅰ）证明见解析；（Ⅱ）$60^\circ$。
  \item[\textbf{20.}] （Ⅰ）$a_n=\dfrac{n^2}{2^{n-1}}$；（Ⅱ）$S_n=5-\dfrac{2n+5}{2^n}$；（Ⅲ）$T_n=12-\dfrac{n^2+4n+6}{2^{n-1}}$。
  \item[\textbf{21.}] （Ⅰ）$C_1:\dfrac{x^2}{36}+\dfrac{y^2}{27}=1$，$C_2:y^2=12x$；（Ⅱ）$|MN|=12$。
  \item[\textbf{22.}] （Ⅰ）在 $(-2,1)$ 上递增，在 $(-\infty,-2)$、$(1,+\infty)$ 上递减，极小值 $f(-2)=-\dfrac12$、极大值 $f(1)=1$；（Ⅱ）$a-b$ 的最大值为 $5$。
\end{itemize}
